% End-evaluation report: present, future and nutrition results (numbers from generated macros)
\section{Which crops are most vulnerable today?}\label{sec:present}

\begin{figure}[t]
\centering
\includegraphics[width=\linewidth]{G01_crop_ranking.pdf}
\caption{Nine-crop vulnerability, 2001--2024 (\Fd). Left: mean over each crop's countries; right:
weighted by each country's production. Colour = crop group.}\label{fig:ranking}
\end{figure}

\paragraph{Ranking.} With equal country weights the order is \orderEq; weighted by production it is
\orderPw\ (Figure~\ref{fig:ranking}, Table~\ref{tab:cropindex}). Wheat is first under both weightings,
in \robTopFfour\ of \robNFfour\ one-at-a-time robustness rankings of \Fd\ and in \mcTopShare\% of 2,000
random weightings (Section~\ref{sec:checks}). Positions below the first are less stable: \secondEq\ is
second with equal weights, \secondPw\ with production weights. By group, \groupOrder\
(Table~\ref{tab:groups}); the group mean describes the three selected crops, not every cereal, pulse or oilseed.

\paragraph{Where vulnerability sits.} The highest pairs are \topPairA\ (\topPairAValue), \topPairB\
(\topPairBValue), \topPairC\ (\topPairCValue), \topPairD\ (\topPairDValue), \topPairE\ (\topPairEValue)
and \topPairF\ (\topPairFValue) (Figure~\ref{fig:matrix}). They share one pattern: a crop that is a large
part of the national diet \emph{and} grown on rainfed land where its water need is often unmet --- wheat in
the Russian, US and Indian plains, maize in South Africa, beans in Brazil and Tanzania. In \nTopSevenWater\ of
the seven highest pairs the water term is larger than the heat term, and over all pairs water accounts for
\presentWaterShare\% of the summed score today; heat matters most for the future increase
(Section~\ref{sec:future}). Wheat in India has the highest wheat heat exposure, but even there the water
term is larger once heat is scaled by the highest heat level in the panel (Table~\ref{tab:worked}).

\paragraph{What drives the scores.} Wheat combines high importance (in the 2021--23 food balance sheets it is
the largest single food source of zinc in \wheatTopZinc, of iron in \wheatTopIron\ and of protein in \wheatTopProtein\ of the \nDivUnits\ units (panel countries and World); \texttt{food\_supply\_diversity.csv}) with exposure on both hazards (Figure~\ref{fig:components}).
Maize is second with equal weights mainly through South Africa (maize supplies \maizeSAkcal\% of national
energy, \maizeSArainfed\% of its area is rainfed, $E_{\text{water}}=\maizeSAEwater$) and Ukraine (the largest water exposure of any
maize pair, \maizeUAEwater). With production weights it drops, because in the two largest producers most maize is not
eaten: only \maizeFoodUSA\% (USA) and \maizeFoodChina\% (China) of domestic supply goes to food (\texttt{crop\_use\_shares.csv}), so
its importance there is small. Rice scores low because
most rice is irrigated (zero water stress in our measure) and because monthly mean temperatures rarely
exceed its \SI{35}{\celsius} anthesis threshold: this is a property of the measurement, not proof that
rice is safe. Pulses enter the top half through common bean (heat) and chickpea (water), but their
national importance is small except in East Africa and South Asia, which keeps lentil low.

\paragraph{How adding crops changed the numbers.} The physical stress of every original pair is unchanged,
but the normalisation reference is not: the largest heat level is now \refHeatPair\ (\refHeat\,\si{\celsius}
above its threshold), and the largest water level \refWaterPair\ (\refWater). For wheat in India the raw
heat level is \exNineLheatc\,\si{\celsius} in both runs, but its normalised exposure falls from \exFourEheat\
to \exNineEheat\ (Table~\ref{tab:worked}), so all original scores roughly halve while their order is
unchanged (four-crop order \orderOrigFourEq; within the nine-crop run \orderOrigNineEq). The bean reference
is itself uncertain (a chamber-derived threshold from a heat-susceptible Andean bean applied to Indian
``beans'', many of which are tropical \emph{Vigna} pulses); raising bean's threshold by \SI{2}{\celsius}
or all new legumes' by \SI{3}{\celsius} leaves wheat first and bean third (Table~\ref{tab:robust}).

% worked example: wheat in India, from raw components to F4 (values from worked_example.csv via macros)
\begin{table}[t]
\centering\small
\caption{Worked example: wheat in India. The physical inputs are identical in both runs; only the
normalisation maxima differ. Values from \texttt{worked\_example.csv}; the hand calculation
$\tfrac12(E_TS_T+E_PS_P)\tilde I$ equals the pipeline value (checked by an assertion in the code).}\label{tab:worked}
\begin{tabular}{lrr}
\toprule
 & four-crop run & nine-crop run\\
\midrule
Heat level $L_{\text{heat}}$ (\si{\celsius} of $\Teff$ above \SI{25}{\celsius}) & \exFourLheatc & \exNineLheatc\\
\quad maximum over pairs & \exFourmaxLheat & \exNinemaxLheat\\
\quad $E_{\text{heat}}$ & \exFourEheat & \exNineEheat\\
Water level $L_{\text{water}}$ (rainfed share $\times$ deficit) & \exFourLwater & \exNineLwater\\
\quad maximum over pairs & \exFourmaxLwater & \exNinemaxLwater\\
\quad $E_{\text{water}}$ & \exFourEwater & \exNineEwater\\
Heat slope $s_T$ (/\si{\celsius}) and maximum over crops & \exFoursTpc\ / \exFourmaxsT & \exNinesTpc\ / \exNinemaxsT\\
\quad $S_T$ & \exFourSheat & \exNineSheat\\
$K_y$ and maximum over crops & \exFourKy\ / \exFourmaxKy & \exNineKy\ / \exNinemaxKy\\
\quad $S_P$ & \exFourSwater & \exNineSwater\\
Importance $I$ and maximum & \exFourI\ / \exFourmaxI & \exNineI\ / \exNinemaxI\\
\quad $\tilde I$ & \exFourInorm & \exNineInorm\\
\midrule
\Fd & \exFourFFourpipeline & \exNineFFourpipeline\\
\bottomrule
\end{tabular}
\end{table}


\section{How does vulnerability change by 2041--2060?}\label{sec:future}

\begin{figure}[t]
\centering
\includegraphics[width=\linewidth]{G04_present_future_crops.pdf}
\caption{Historical (1971--2000, open) and future (2041--2060, filled) scores, MRI-ESM2-0 SSP2-4.5,
change factors on the same years; one fixed normalisation reference.}\label{fig:future}
\end{figure}

Under one model and one scenario, the flowering-window effective temperature rises by
\meanDTeff\,\si{\celsius} on average over the \nPairs\ pairs. The score rises in \nPairsUp\ of \nPairs\ pairs;
\nPairsExceedRef\ pair exceeds the present-day reference maximum. The crop order changes little:
\futOrderEq\ (equal weights; historical: \histOrderEq). The largest increase is \futTopPair\
(\futTopPairChange), and summed over all pairs the heat term contributes about \futHeatShare\% of the increase and the water term \futWaterShare\%
(Table~\ref{tab:futurepairs}). The
climatology check (thresholds applied to the 20-year means) gives the same sign of change in
\climCheckSameSign\% of pairs (Spearman $\rho$ = \climCheckRho), so the result does not hinge on the
change-factor choice. Every pair whose score falls (rice, maize, soybean and groundnut in some countries, e.g.\ maize in South
Africa and rice in Nigeria) falls through the water term alone: the change factors give a smaller rainfed
deficit ($1-P/\ETc$) in their growing season, while their heat term stays near zero. A crop can stay nutritionally important while its climate stress rises; and a low score can
still reflect the monthly heat measure.

\begin{tabular}{lrrrrrrcrrc}
\toprule
Crop & $\Delta T_{eff}$ & 2001--24 & 1971--2000 & 2041--60 & $\Delta$F4 & \% & rank eq & pairs up & F4 pw 2041--60 & rank pw \\
\midrule
Wheat & +2.0 & 0.0563 & 0.0535 & 0.0616 & +0.0080 & +15\% & 1$\to$1 & 6/6 & 0.0762 & 1$\to$1 \\
Maize & +2.0 & 0.0297 & 0.0263 & 0.0254 & -0.0009 & -3\% & 2$\to$2 & 3/6 & 0.0084 & 6$\to$6 \\
Common bean & +2.0 & 0.0172 & 0.0167 & 0.0224 & +0.0057 & +34\% & 3$\to$3 & 7/7 & 0.0276 & 3$\to$2 \\
Rapeseed-mustard & +2.0 & 0.0078 & 0.0076 & 0.0081 & +0.0006 & +8\% & 5$\to$4 & 7/7 & 0.0161 & 4$\to$4 \\
Soybean & +1.9 & 0.0078 & 0.0079 & 0.0080 & +0.0001 & +1\% & 4$\to$5 & 3/7 & 0.0121 & 5$\to$5 \\
Chickpea & +2.0 & 0.0053 & 0.0052 & 0.0053 & +0.0001 & +2\% & 6$\to$6 & 3/7 & 0.0227 & 2$\to$3 \\
Rice & +1.8 & 0.0024 & 0.0027 & 0.0022 & -0.0005 & -17\% & 7$\to$7 & 5/7 & 0.0016 & 9$\to$9 \\
Lentil & +2.2 & 0.0021 & 0.0021 & 0.0021 & -0.0000 & -1\% & 9$\to$8 & 4/7 & 0.0040 & 7$\to$7 \\
Groundnut & +1.6 & 0.0015 & 0.0021 & 0.0017 & -0.0004 & -19\% & 8$\to$9 & 2/5 & 0.0019 & 8$\to$8 \\
\bottomrule
\end{tabular}
\par\vspace{2pt}{\footnotesize MRI-ESM2-0, SSP2-4.5, change factors on 1971--2000 years; one fixed normalisation reference (2001--24 nine-crop maxima); sensitivity, importance, crop area, irrigation and calendars held fixed. $\Delta T_{eff}$ = mean change in flowering-window effective temperature (\si{\celsius}). \% only where the historical score is $\geq$1\% of the largest.}


\section{What does a crop loss mean for nutrient supply?}\label{sec:nutrition}

\paragraph{What each crop supplies.} Cereals dominate energy, protein, zinc and iron; the oilseeds mainly
supply fat (as oil); pulses matter for protein and iron mostly in South Asia and East Africa
(Figure~\ref{fig:profiles}).

\paragraph{Adequacy.} Comparing national supply with WHO/FAO recommended intakes, iron is the nutrient
most often below requirement (\nIronShort\ of \nAdqCountries\ panel countries) and zinc never is. This
claim depends on the requirement used: with iron at 15\% bioavailability \nIronShortAlt\ countries fall
short. Recommended intakes are upper-percentile individual requirements compared with a national
\emph{mean} supply, so a national ratio above one does not mean no one is deficient, and a ratio below
one does not by itself mean deficiency (Table~\ref{tab:adequacy}).

\paragraph{A standard 10\% yield loss.} Removing 10\% of a crop's production share of supply (no market
response) changes national MAR by less than 0.005 points in \shockZeroMar\ of \shockN\ pairs, because
national supply stays above the requirement for most nutrients; the largest MAR fall is \shockTopMar\
(\shockTopMarValue\ points) and the largest energy loss \shockTopKcal\ (\shockTopKcalValue\% of national
kcal). Because MAR caps each ratio at one, it hides losses that still leave supply above the requirement;
the nutrient losses themselves are therefore reported alongside (Table~\ref{tab:shock}).

\paragraph{History.} In the old food balance sheets (1961--2013), a bad harvest significantly lowered food
supply in \passOldPos\ of \passOldN\ pairs; in the new series (2010--2023) in \passNewPos\ of \passNewN.
These are associations: trade and stocks buffer most losses, and insignificant estimates are not proof of
buffering.

\paragraph{Future yield and nutrition (conditional).} Pairing physical stress with the literature response
(heat: $s_T$ per \si{\celsius} above the threshold; water: $K_y$ times the rainfed deficit) gives the largest
climate-driven yield change for \futYieldWorst\ (mean \futYieldWorstValue\%), and the largest nutrient
consequence for \futNutTopPair\ (MAR \futNutTopMar\ points) (Table~\ref{tab:futyield}). The empirical
regressions do not support these numbers pair by pair (rank correlation \corrLitEmp): most empirical
heat coefficients are not statistically detectable, and the two approaches answer different questions.
We therefore treat the future nutrient numbers as a conditional screening estimate, not a forecast.
